% New contribution relative to the inspected 2026-11-10 draft corpus.
% This fragment is intended for inclusion after standard theorem environments.
\section{Conjunctive induction still leaves a symmetry obstruction}
\label{sec:pp-obstruction}

The earlier eventually-constant-power obstruction in the inspected
\texttt{ipra\_research\_2026-11-10} draft, specifically
\texttt{theory/finite\_parameter\_obstruction.tex}, includes equation
induction, finite conjunctions of equations, and universal Horn
arithmetic laws~\cite{research}.  That source already attributes the finite-parameter
involution argument to earlier unary-amalgam work; neither the power
construction nor the automorphism method is claimed here as new.
The result below first extends that particular construction to existential
witnesses: both the induction matrices and the regular axioms may contain
positive existential quantification.  A further closure theorem allows
unrestricted finite alternation of existential and universal quantifiers,
provided that the formulas use conjunction and no disjunction,
implication, or negation internally.  Both layers extend the existing
construction; neither is presented as a literature-wide priority claim
or a positive axiomatization result.

\subsection{Exact hypotheses}

Let $\mathcal N$ be any one-sorted finitary expansion of the standard
structure $(\mathbb N,0,S)$ by total functions and relations.  The
signature can be infinite; each symbol has finite arity, and each formula
is finite.  In particular, it can contain every primitive-recursive
function.  A \emph{primitive-positive} formula is built from atomic
formulas and $\top$ by finite conjunction and existential quantification.
Neither disjunction, implication, negation, nor universal quantification
is permitted inside such a formula.

We permit two kinds of regular axioms:
\[
 \forall\vec x\,\bigl(\varphi(\vec x)\to\psi(\vec x)\bigr),
 \qquad
 \forall\vec x\,\bigl(\varphi(\vec x)\to\bot\bigr),
 \tag{Reg}\label{eq:pp-regular}
\]
where $\varphi,\psi$ are primitive-positive and the displayed sentence
is true in $\mathcal N$.  The second form explicitly includes negative
constraints such as $0=S0\to\bot$.  Terminology differs on whether
regular logic includes falsity, so this form is listed separately.

For every primitive-positive $A(n,\vec p)$, include the full induction
sentence
\[
 \forall\vec p\left[
 \left(A(0,\vec p)\wedge
       \forall n\,(A(n,\vec p)\to A(Sn,\vec p))\right)
 \to \forall n\,A(n,\vec p)\right].
 \tag{PPI}\label{eq:pp-induction}
\]
Parameters $\vec p$ remain fixed during induction.  This is equivalent
to the usual contextual induction rule for this class of matrices in
intuitionistic first-order logic.  Retained assumptions can contain
arbitrary formulas.  The claim below does not restrict them to
primitive-positive formulas.

Let $T_{\mathrm{pp}}(\mathcal N)$ be intuitionistic first-order logic
with equality, all sentences \eqref{eq:pp-regular}, and all instances
\eqref{eq:pp-induction}.  This semantic upper theory need not have a
recursive axiom set.  The obstruction also applies to any recursively
axiomatized subtheory containing the same symbols and using only axioms
from these classes.  More generally, it applies to any theory satisfied
by the model constructed below.

\subsection{The model and its positive existential formulas}

Write
\[
 E=\{a\colon\mathbb N\to\mathbb N:
       a\text{ is eventually constant}\}.
\]
Interpret every function of $\mathcal N$ coordinatewise.  For an
$r$-ary relation $R$, put
\[
 E\models R(a_1,\ldots,a_r)
 \quad\Longleftrightarrow\quad
 \mathcal N\models R(a_1(i),\ldots,a_r(i))
       \text{ for every }i\in\mathbb N.
\]
Equality is actual equality of sequences, and constants are constant
sequences.  Finitarity ensures that coordinatewise functions take
eventually constant inputs to eventually constant outputs.  Denote the
resulting structure by $\mathcal E$.

\begin{lemma}[Primitive-positive transfer]
\label{lem:pp-transfer}
For every primitive-positive $\varphi(\vec x)$ and every finite tuple
$\vec a$ from $E$,
\[
 \mathcal E\models\varphi(\vec a)
 \quad\Longleftrightarrow\quad
 \mathcal N\models\varphi(\vec a(i))
                 \text{ for every }i\in\mathbb N.
 \tag{PP}\label{eq:pp-transfer}
\]
\end{lemma}
\begin{proof}
Use external induction on the fixed finite formula.  Atomic formulas
and finite conjunctions follow immediately from the coordinatewise
definitions.  The forward existential direction follows by evaluating
a witness at each coordinate.

For the reverse existential direction, suppose
$\mathcal N\models\exists y\,\theta(y,\vec a(i))$ at every
coordinate, and assume the transfer statement for $\theta$.  There is
a finite set $F$ such that $\vec a(i)=\vec a_\infty$ for every
$i\notin F$.  Choose a witness $b_\infty$ for the common tail tuple,
and a witness $b_i$ separately for each of the finitely many coordinates
$i\in F$.  The sequence with values $b_i$ on $F$ and $b_\infty$
elsewhere belongs to $E$.  The induction hypothesis gives
$\mathcal E\models\theta(b,\vec a)$, hence the required
existential statement.  Only finitely many witness choices are used.
Nested or tuple existential quantifiers are handled by repeated use of
this argument.
\end{proof}

\begin{lemma}[Regular axioms and primitive-positive induction]
\label{lem:pp-model}
The classical structure $\mathcal E$ satisfies every axiom of
$T_{\mathrm{pp}}(\mathcal N)$.
\end{lemma}
\begin{proof}
For a regular axiom, a true antecedent in $\mathcal E$ gives a true
antecedent at every coordinate by Lemma~\ref{lem:pp-transfer}.  Validity
in $\mathcal N$ gives the conclusion at every coordinate; the reverse
direction of the lemma reconstructs its witnesses in $\mathcal E$.
For a negative regular constraint, a true antecedent in $\mathcal E$
would give a true antecedent in $\mathcal N$ at coordinate zero,
contrary to the constraint's validity there.

Fix an instance of \eqref{eq:pp-induction} and a parameter tuple
$\vec p\in E$.  Suppose its base and progress hypotheses hold in
$\mathcal E$.  Fix a coordinate $i$.  The transfer lemma gives
$\mathcal N\models A(0,\vec p(i))$.  If
$\mathcal N\models A(k,\vec p(i))$, form the sequence $a$
having value $k$ at coordinate $i$ and value $0$ elsewhere.  At
coordinate $i$, $A(a(i),\vec p(i))$ holds by the supposition;
at every other coordinate it holds by the transferred base hypothesis.
The reverse transfer lemma therefore gives
$\mathcal E\models A(a,\vec p)$.  Apply the assumed progress
sentence in $\mathcal E$, and then transfer back at coordinate $i$.
The result is
$\mathcal N\models A(k+1,\vec p(i))$.

Ordinary external induction on the standard number $k$ now gives
$\mathcal N\models A(k,\vec p(i))$ for every $k$.  This holds
at each coordinate $i$.  For any $a\in E$, the reverse transfer
lemma gives $\mathcal E\models A(a,\vec p)$.  Hence the whole
induction instance holds.  The parameter tuple has remained fixed
throughout.  Nothing here assumes that witnesses for $A$ are terms or
are uniformly definable functions of $n$.
\end{proof}

A classical model is a one-node intuitionistic Kripke model.  Thus the
lemma supplies a model of the intuitionistic theory, not merely evidence
about its classical restriction.  In particular, any derivation using
arbitrarily quantified retained assumptions must be valid in this model.

For orientation, successor coverage already fails on the original
domain.  The sequence equal to $1$ at one coordinate and $0$ elsewhere
is not the zero sequence and is not a successor, since every successor
sequence has a positive value at every coordinate.  This observation
alone would not exclude an interpretation on a different definable
domain; the following symmetry argument is what handles that freedom.

\subsection{The automorphism obstruction}

\begin{lemma}[Finite parameters leave an involution]
\label{lem:pp-involution}
For every finite tuple $\vec p\in E$, the structure
$(\mathcal E,\vec p)$ has a nonidentity automorphism of order two.
\end{lemma}
\begin{proof}
Outside a finite set, every component of $\vec p$ is constant.  Choose
two distinct coordinates $i,j$ outside this set and swap them.
Coordinatewise interpretation makes the swap an automorphism preserving
all function and relation symbols, irrespective of the size of the
signature.  It fixes $\vec p$ and has order two.  It is nonidentity
because the sequence equal to $1$ at $i$ and $0$ elsewhere is moved to
the corresponding sequence supported at $j$.
\end{proof}

\begin{theorem}[No ordinary bi-interpretation with iPRA]
\label{thm:pp-obstruction}
$T_{\mathrm{pp}}(\mathcal N)$ is not bi-interpretable with iPRA by
ordinary finitary first-order syntactic interpretations, of finite
dimension, whose composites are provably definably isomorphic to the
respective identity interpretations.  Arbitrary formula-defined domains,
relation/function graphs, and quotient equality are allowed.  Each
model may use a fixed finite parameter tuple, including parameters for
the return isomorphisms, as long as the claimed interpretations and
isomorphisms are required to exist in that model.  In particular, this
covers a syntactic parameter scheme whose finite parameter condition is
provably inhabited and under which all these facts are proved uniformly.
The same conclusion holds for every theory having $\mathcal E$ as a
model.  A parameter condition that simply excludes $\mathcal E$ and is
not proved inhabited there is outside the claim.
\end{theorem}
\begin{proof}
Every classical model of iPRA has a definable linear order.  For example,
in the standard primitive-recursive language use
$x\le y\leftrightarrow\exists z\,(x+z=y)$; iPRA proves that this
is a linear order.  An order-preserving automorphism of finite order is
the identity: if $a<\sigma(a)$, iterating the strict inequality gives
$a<\sigma(a)<\cdots<\sigma^r(a)=a$, and the case
$\sigma(a)<a$ is analogous.

Suppose there were a bi-interpretation.  Apply the interpretation to
$\mathcal E$ to obtain a classical iPRA model $\mathcal M$.
Collect the finitely many parameters of both interpretations and their
definable return isomorphisms; representatives of any parameters in an
interpreted quotient add only finitely many elements of $E$.
Lemma~\ref{lem:pp-involution} supplies a nonidentity involution
$\sigma$ fixing all of them.

Every first-order interpretation carries such an automorphism to the
interpreted structure, including when its domain is a definable set of
tuples modulo a definable equivalence relation.  The induced
automorphism $\widehat\sigma$ of $\mathcal M$ has order dividing
two, so the preceding linear-order argument makes it the identity.
Its further interpretation is therefore the identity on the
double-interpreted structure.  The source-side return isomorphism is
definable with parameters fixed by $\sigma$, hence is equivariant:
\[
 \eta(\sigma(a))=\widehat{\widehat\sigma}(\eta(a))=\eta(a).
\]
Since $\eta$ is injective, $\sigma(a)=a$ for every $a\in E$,
contrary to the choice of $\sigma$.
\end{proof}

\begin{corollary}[Additional sorts do not evade this obstruction]
\label{cor:pp-many-sorted}
Let $\mathcal A$ be a finitary many-sorted structure with nonempty
sorts and a distinguished numerical sort that is the standard
$(\mathbb N,0,S)$.  Other sorts may, for example, consist of finitary
objects or specified numerical functions, with arbitrary total finitary
operations and relations between the sorts.  Form the intuitionistic
theory of all true many-sorted regular sequents and negative regular
constraints, together with induction on the numerical sort for every
primitive-positive matrix.  It is not bi-interpretable with iPRA by
the ordinary syntactic interpretations of
Theorem~\ref{thm:pp-obstruction}.  Interpretation domains may use finite
products of any of the source sorts and definable quotients.
\end{corollary}
\begin{proof}
For each sort with domain $A_s$, take the eventually constant sequences
$\mathbb N\to A_s$, and interpret all symbols coordinatewise, using
the same coordinate index set for every sort.  A fixed formula and its
parameter tuple involve only finitely many sorts and coordinates with
exceptional values.  The primitive-positive transfer proof therefore
works unchanged: existential witnesses have their specified sorts,
and finite exceptional witnesses plus one common tail witness give
an eventually constant witness of that sort.  This proves every true
regular axiom.

For induction, the varying sequence $a_{i,k}$ belongs to the
distinguished numerical sort.  All other parameter sorts stay fixed.
The proof of Lemma~\ref{lem:pp-model} consequently proves every
primitive-positive numerical-induction instance, even when witnesses
and parameters range over other sorts.

Swap two coordinates outside the combined finite support of all
interpretation and return-map parameters, simultaneously in every
sort.  This is an automorphism of the whole many-sorted structure.
It is nonidentity already on the numerical sort, by the sequence
supported at one of the swapped coordinates.  Its induced
automorphism on an interpreted iPRA model must be the identity.
Definability of the source return isomorphism, on every source sort,
then makes the original automorphism the identity on every sort.
The numerical-sort contradiction suffices.  No faithfulness of the
action on each other sort is assumed or needed.
\end{proof}

\subsection{Universal quantifiers give a stronger closure theorem}

The primitive-positive result is not the limit of this method.  Let
$\mathcal C$ be the class of \emph{conjunctive formulas} generated from
atomic formulas and $\top$ by $\wedge$, $\exists$, and $\forall$.
Quantifier alternation is unrestricted.  There is still no disjunction,
implication, or negation inside a matrix from $\mathcal C$.
The explicit grammar is the convention used here; no identification
with a differently defined syntactic fragment is assumed.

\begin{theorem}[Conjunctive transfer and induction]
\label{thm:conjunctive-induction}
For every $\varphi(\vec x)\in\mathcal C$,
\[
 \mathcal E\models\varphi(\vec a)
 \quad\Longleftrightarrow\quad
 \mathcal N\models\varphi(\vec a(i))
                  \text{ for every }i\in\mathbb N.
\]
Consequently $\mathcal E$ satisfies induction for every matrix in
$\mathcal C$, as well as every sentence true in $\mathcal N$ of the
forms
\[
 \forall\vec x\,\bigl(\varphi(\vec x)\to\psi(\vec x)\bigr),
 \qquad
 \forall\vec x\,\bigl(\varphi(\vec x)\to\bot\bigr),
 \qquad \varphi,\psi\in\mathcal C.
\]
These conclusions hold also in the many-sorted setting of
Corollary~\ref{cor:pp-many-sorted}, with induction in the numerical
sort.  The resulting stronger theories are still not bi-interpretable
with iPRA under the interpretations specified in
Theorem~\ref{thm:pp-obstruction}.
\end{theorem}
\begin{proof}
Induct on the fixed finite formula.  The atomic and conjunction cases
are unchanged.  The existential proof of
Lemma~\ref{lem:pp-transfer} uses only that the body already satisfies
transfer and that its finite parameter tuple has a common tail.
It therefore works even when the body contains universal quantifiers.

For the universal case, suppose
$\mathcal E\models\forall y\,\theta(y,\vec a)$.  Fix a
coordinate $i$ and any element $b$ of the corresponding standard sort.
The constant sequence $\overline b$ is an element of the sequence
sort.  Instantiate the universal statement at $\overline b$ and apply
the induction hypothesis at coordinate $i$.  This yields
$\mathcal N\models\theta(b,\vec a(i))$.  As $b$ was arbitrary,
the coordinate satisfies the universal statement.  Conversely, if
every coordinate satisfies that universal statement, then for every
sequence $b$ every coordinate satisfies $\theta(b(i),\vec a(i))$.
The induction hypothesis gives
$\mathcal E\models\theta(b,\vec a)$ for every such $b$,
which is the required universal statement in $\mathcal E$.

The proofs of regular-axiom preservation and induction in
Lemma~\ref{lem:pp-model} use only exact coordinatewise transfer;
apply them to this larger class.  The finite-parameter involution and
return-isomorphism obstruction use only the unchanged underlying model,
so they apply verbatim.  The proof allows arbitrarily alternating
quantifiers; it supplies no term-witness or decidability assumption.
\end{proof}

The one-sorted primitive-positive and many-sorted statements above are
thus special cases.  Even unbounded quantifier alternation in this
conjunctive fragment does not force the uniform branch information
missing from successor coverage.

\paragraph{Function objects and extensionality.}
The many-sorted statement also applies when a sort consists of specified
actual numerical functions and the signature includes a total finitary
application operation. Extensionality has the allowed form
\[
 \bigl(\forall n\;\operatorname{app}(f,n)=\operatorname{app}(g,n)\bigr)
       \to f=g.
\]
Its antecedent is conjunctive, so the sequence model preserves it.
Thus adding function objects, application, and extensionality does not
by itself evade the obstruction. Equational combinator laws and any
other standard-true conjunctive sequents may be included as well.
This observation does not address a calculus with additional principles
outside the stated class; it identifies the exact assumptions under
which that possible change of domain still retains the symmetry.

\paragraph{Scope.}
The theorem is about ordinary first-order bi-interpretations, not
equiconsistency alone.  It allows existentially specified witnesses and
conjunctive induction of arbitrary quantifier alternation, so changing
a term-defined operation to a positive existential graph does not by
itself evade this obstruction.
It does not rule out arbitrary one-way interpretations, nonregular
logical rules, disjunctive coverage axioms, or arbitrary formula-defined
expansions that are not interpreted pointwise.  Nor does it yield a
numerical lower bound for proof-checker source size.  Its practical use
is to rule out an entire family of proposed positive existential
completions before attempting a full bi-interpretation proof.
